\documentclass[12pt]{article}
\usepackage[utf8]{inputenc}
\usepackage[spanish]{babel}
\usepackage{amsmath}
\usepackage{amsfonts}
\usepackage{amssymb}
\usepackage{geometry}
\geometry{letterpaper, margin=2.5cm}

\begin{document}

\begin{center}
    \Large \textbf{Evaluación de Examen 1: Unidad I} \\
    \large Física para Ciencias de la Salud (005-1713) \\
    \vspace{0.5cm}
    \textbf{Estudiante:} Thomairys Sosa \quad \textbf{C.I.:} 33.685.285
    \vspace{0.3cm} \\
    \large \textbf{Calificación Final: 8/10 puntos}
\end{center}

\vspace{0.5cm}

\section*{Análisis de Resultados}

\subsection*{1. Ángulo plano (Conversión a radianes)}
\textbf{Procedimiento y resultado:} Plantea correctamente la conversión usando el factor $\frac{\pi}{180^\circ}$. Simplifica la fracción a $\frac{5\pi}{12}$ y además calcula su valor decimal ($1,309$ rad). Procedimiento completo y correcto. \\
\textbf{Veredicto:} \textbf{Correcto} (2/2 pts).

\subsection*{2. Sistemas de coordenadas (Longitud del hueso)}
\textbf{Procedimiento y resultado:} Excelente resolución. Escribe la fórmula de la distancia, identifica las coordenadas y calcula las diferencias en los ejes ($x=6, y=8$). Aplica el Teorema de Pitágoras y llega directamente a la raíz de 100, obteniendo el resultado exacto de $10$ cm. \\
\textbf{Veredicto:} \textbf{Correcto} (2/2 pts).

\subsection*{3. Vectores (Fuerza resultante)}
\textbf{Procedimiento y resultado:} Aunque la agrupación inicial de términos es un poco confusa en la escritura, el resultado final es perfecto. Suma correctamente las componentes ($45-15=30$ para $\hat{i}$ y $25+35=60$ para $\hat{j}$), obteniendo el vector resultante exacto $(30\hat{i} + 60\hat{j})$ N. \\
\textbf{Veredicto:} \textbf{Correcto} (2/2 pts).

\subsection*{4. Potencias de diez (Notación científica y prefijo)}
\textbf{Procedimiento y resultado:} Transforma la cifra decimal correctamente a una expresión con potencias de base diez ($12,5 \times 10^{-6}$). Sin embargo, omite la segunda instrucción de la pregunta: identificar y nombrar explícitamente el prefijo del S.I. (micrómetros). \\
\textbf{Veredicto:} \textbf{Incompleto / Parcialmente correcto} (1/2 pts).

\subsection*{5. Sistemas de unidades (Conversión de densidad)}
\textbf{Procedimiento y resultado:} Intenta resolver el problema aplicando factores de conversión en cadena. Plantea correctamente el factor de gramos a kilogramos ($1 \text{ kg} / 1000 \text{ g}$), pero comete un \textbf{error conceptual grave en la conversión de volumen}, escribiendo que $1 \text{ m}^3 = 100 \text{ cm}^3$ en lugar del valor correcto ($1.000.000 \text{ cm}^3$). Por causa de este factor incorrecto, obtiene un resultado erróneo de $0,103 \text{ kg/m}^3$. \\
\textbf{Veredicto:} \textbf{Parcialmente correcto} (Estructura de factores de conversión planteada, pero error en equivalencia cúbica) (1/2 pts).

\vspace{0.5cm}
\section*{Comentarios Generales y Sugerencias}
La estudiante demuestra un buen dominio en trigonometría, manejo de fórmulas geométricas y sumas de álgebra vectorial.
\begin{itemize}
    \item Se recomienda repasar fuertemente las conversiones de unidades al cubo (volumen), recordando que al pasar de una unidad lineal (cm) a otra cúbica ($\text{cm}^3$), el factor numérico también debe elevarse al cubo (Ejemplo: $100^3 = 1.000.000$).
\end{itemize}

\end{document}